\import{algebra/module/direct-sum.tex}
\import{algebra/module/linear-map.tex}
\import{algebra/ring/regular-module.tex}

\section{Free modules}

The free module on an index set is the direct sum of one copy of the left
regular module at each index.  Thus its vectors are the existing finitely
supported sections, its basis vectors are the existing summand inclusions,
and its universal map is the existing direct-sum mediator.

\begin{definition}\label{free_module}
    $\FreeModule{R}{I}
        =\DirectSumModule{R}{I}{I\times\{\LeftRegularModule{R}\}}$.
\end{definition}

\begin{proposition}\label{free_module_constructor_eq}
    $\FreeModule{R}{I}
        =\DirectSumModule{R}{I}{I\times\{\LeftRegularModule{R}\}}$.
\end{proposition}
\begin{proof}
    Follows by \obv{free_module}.
\end{proof}


\begin{proposition}\label{free_module_regular_family_type}
    $I\times\{\LeftRegularModule{R}\}\in
        \funs{I}{\{\LeftRegularModule{R}\}}$.
\end{proposition}
\begin{proof}
    Follows by \obv{constant_graph_type}.
\end{proof}

\begin{proposition}\label{free_module_regular_family_apply}
    Let $i\in I$.
    Then $(I\times\{\LeftRegularModule{R}\})(i)
        =\LeftRegularModule{R}$.
\end{proposition}
\begin{proof}
    Follows by \obv{constant_graph_application}.
\end{proof}

\begin{theorem}\label{free_module_is_exact}
    Let $R$ be a ring.
    Then $\FreeModule{R}{I}$ is an exact left module over $R$.
\end{theorem}
\begin{proof}
    Let $F=I\times\{\LeftRegularModule{R}\}$.
    $F$ is a function on $I$ by
        \auto{free_module_regular_family_type,
            funs_is_function,funs}.
    For every $i\in I$ we have $F(i)$ is a left module over $R$ by
        \auto{free_module_regular_family_apply,
            left_regular_module_is_exact}.
    Follows by
        \auto{free_module_constructor_eq,direct_sum_module_is_exact}.
\end{proof}

\begin{proposition}\label{free_module_carrier_family}
    Let $R$ be a ring.
    Let $F=I\times\{\LeftRegularModule{R}\}$ and
        $A=\moduleAdditiveGroupFamily{I}{F}$.
    Then
        $\productMonoidCarrierFamily{I}{A}
            =I\times\{\carrier[R]\}$.
\end{proposition}
\begin{proof}
    Let $C=\productMonoidCarrierFamily{I}{A}$ and
        $K=I\times\{\carrier[R]\}$.
    $F$ is a function on $I$ by
        \auto{free_module_regular_family_type,
            funs_is_function,funs}.
    Hence $A$ is a function on $I$ by
        \auto{module_additive_group_family_type}.
    Hence $C$ is a function and $\dom{C}=I$ by
        \auto{product_monoid_carrier_family_type,
            funs_is_function,funs}.
    $K\in\funs{I}{\{\carrier[R]\}}$ by
        \auto{constant_graph_type}.
    Hence $K$ is a function and $\dom{K}=I$ by
        \obv{funs_is_function,funs}.
    Show that for every $i\in I$ we have $C(i)=K(i)$.
    \begin{subproof}
        Fix $i\in I$.
        $C(i)=\carrier[A(i)]$ by
            \auto{product_monoid_carrier_family_apply}.
        $A(i)=\ModuleAdditiveGroup{F(i)}$ by
            \auto{module_additive_group_family_apply}.
        $F(i)=\LeftRegularModule{R}$ by
            \auto{free_module_regular_family_apply}.
        $\ModuleAdditiveGroup{\LeftRegularModule{R}}
            =\AdditiveGroup{R}$ by
            \obv{left_regular_module_additive_group}.
        Hence $\carrier[A(i)]=\carrier[R]$ by
            \auto{additive_group_fields}.
        $K(i)=\carrier[R]$ by \auto{constant_graph_application}.
    \end{subproof}
    Follows by \auto{funext_on_domain}.
\end{proof}

\begin{proposition}\label{free_module_zero_section}
    Let $R$ be a ring.
    Let $F=I\times\{\LeftRegularModule{R}\}$ and
        $A=\moduleAdditiveGroupFamily{I}{F}$.
    Then
        $\productMonoidNeutral{I}{A}
            =I\times\{\ringZero[R]\}$.
\end{proposition}
\begin{proof}
    Let $z=\productMonoidNeutral{I}{A}$ and
        $e=I\times\{\ringZero[R]\}$.
    $F$ is a function on $I$ by
        \auto{free_module_regular_family_type,
            funs_is_function,funs}.
    Hence $A$ is a function on $I$ by
        \auto{module_additive_group_family_type}.
    For every $i\in I$ we have $A(i)$ is a commutative group by
        \auto{module_additive_group_family_apply,
            free_module_regular_family_apply,
            left_regular_module_is_exact,
            module_additive_group_is_exact_commutative_group}.
    $z\in\dependentProduct{I}{
        \productMonoidCarrierFamily{I}{A}}$ by
        \auto{product_monoid_neutral_type,
            commutative_group_is_group,group_is_monoid}.
    Hence $z$ is a function and $\dom{z}=I$ by
        \auto{dependent_product_iff,funs_is_function,funs}.
    $e\in\funs{I}{\{\ringZero[R]\}}$ by
        \auto{constant_graph_type}.
    Hence $e$ is a function and $\dom{e}=I$ by
        \obv{funs_is_function,funs}.
    Show that for every $i\in I$ we have $z(i)=e(i)$.
    \begin{subproof}
        Fix $i\in I$.
        $z(i)=\neutral[A(i)]$ by
            \auto{product_monoid_neutral_coordinate,
                commutative_group_is_group,group_is_monoid}.
        $A(i)=\ModuleAdditiveGroup{F(i)}$ by
            \auto{module_additive_group_family_apply}.
        $F(i)=\LeftRegularModule{R}$ by
            \auto{free_module_regular_family_apply}.
        $\ModuleAdditiveGroup{\LeftRegularModule{R}}
            =\AdditiveGroup{R}$ by
            \obv{left_regular_module_additive_group}.
        Hence $z(i)=\ringZero[R]$ by
            \auto{additive_group_fields}.
        $e(i)=\ringZero[R]$ by \auto{constant_graph_application}.
    \end{subproof}
    Follows by \auto{funext_on_domain}.
\end{proof}

\begin{theorem}\label{free_module_carrier}
    Let $R$ be a ring.
    Then
        $\carrier[\FreeModule{R}{I}]
            =\finitelySupportedFunctions
                {I}{\carrier[R]}{\ringZero[R]}$.
\end{theorem}
\begin{proof}
    Let $F=I\times\{\LeftRegularModule{R}\}$ and
        $A=\moduleAdditiveGroupFamily{I}{F}$.
    Let $C=\productMonoidCarrierFamily{I}{A}$ and
        $z=\productMonoidNeutral{I}{A}$.
    Let $D=\DirectSumModule{R}{I}{F}$ and
        $Q=\DirectSumGroup{I}{A}$.
    $F$ is a function on $I$ by
        \auto{free_module_regular_family_type,
            funs_is_function,funs}.
    For every $i\in I$ we have $F(i)$ is a left module over $R$ by
        \auto{free_module_regular_family_apply,
            left_regular_module_is_exact}.
    $C=I\times\{\carrier[R]\}$ by
        \obv{free_module_carrier_family}.
    $z=I\times\{\ringZero[R]\}$ by
        \obv{free_module_zero_section}.
    $\carrier[D]=\carrier[Q]$ by
        \auto{direct_sum_module_carrier}.
    $\carrier[Q]
        =\dependentFinitelySupportedFunctions{I}{C}{z}$ by
        \auto{direct_sum_group_carrier}.
    $\ringZero[R]\in\carrier[R]$ by \obv{ring_zero_type}.
    $\dependentFinitelySupportedFunctions
        {I}{I\times\{\carrier[R]\}}{I\times\{\ringZero[R]\}}
        =\finitelySupportedFunctions
            {I}{\carrier[R]}{\ringZero[R]}$ by
        \obv{constant_family_finite_support_space}.
    Follows by \auto{free_module_constructor_eq}.
\end{proof}

\subsection{Canonical basis vectors}

\begin{definition}\label{free_basis_vector}
    $\freeBasisVector{R}{I}{i}
        =\apply{\directSumModuleInclusion
            {R}{I}{I\times\{\LeftRegularModule{R}\}}{i}}{
                \ringOne[R]}$.
\end{definition}

\begin{proposition}\label{free_basis_vector_type}
    Let $R$ be a ring and $i\in I$.
    Then $\freeBasisVector{R}{I}{i}
        \in\carrier[\FreeModule{R}{I}]$.
\end{proposition}
\begin{proof}
    Let $F=I\times\{\LeftRegularModule{R}\}$.
    $F$ is a function on $I$ by
        \auto{free_module_regular_family_type,
            funs_is_function,funs}.
    For every $j\in I$ we have $F(j)$ is a left module over $R$ by
        \auto{free_module_regular_family_apply,
            left_regular_module_is_exact}.
    $\directSumModuleInclusion{R}{I}{F}{i}\in
        \funs{\carrier[F(i)]}{
            \carrier[\DirectSumModule{R}{I}{F}]}$ by
        \obv{direct_sum_module_inclusion_type}.
    $\carrier[F(i)]=\carrier[R]$ by
        \auto{free_module_regular_family_apply,
            left_regular_module_carrier}.
    $\ringOne[R]\in\carrier[F(i)]$ by \obv{ring_one_type}.
    $\apply{\directSumModuleInclusion{R}{I}{F}{i}}{\ringOne[R]}
        \in\carrier[\DirectSumModule{R}{I}{F}]$ by
        \obv{funs_type_apply}.
    Follows by
        \auto{free_basis_vector,free_module_constructor_eq}.
\end{proof}

\begin{proposition}\label{free_basis_vector_point_mass}
    Let $R$ be a ring and $i\in I$.
    Then
        $\freeBasisVector{R}{I}{i}
            =\pointMass{I}{\carrier[R]}{\ringZero[R]}{i}{
                \ringOne[R]}$.
\end{proposition}
\begin{proof}
    Let $F=I\times\{\LeftRegularModule{R}\}$ and
        $A=\moduleAdditiveGroupFamily{I}{F}$.
    Let $C=\productMonoidCarrierFamily{I}{A}$ and
        $z=\productMonoidNeutral{I}{A}$.
    $F$ is a function on $I$ by
        \auto{free_module_regular_family_type,
            funs_is_function,funs}.
    For every $j\in I$ we have $F(j)$ is a left module over $R$ by
        \auto{free_module_regular_family_apply,
            left_regular_module_is_exact}.
    $C=I\times\{\carrier[R]\}$ by
        \obv{free_module_carrier_family}.
    $z=I\times\{\ringZero[R]\}$ by
        \obv{free_module_zero_section}.
    $\apply{\directSumModuleInclusion{R}{I}{F}{i}}{\ringOne[R]}
        =\dependentPointMass{I}{C}{z}{i}{\ringOne[R]}$ by
        \auto{direct_sum_module_inclusion_apply,
            free_module_regular_family_apply,
            left_regular_module_carrier,ring_one_type}.
    Follows by \auto{free_basis_vector,point_mass}.
\end{proof}

\begin{proposition}\label{free_basis_vector_selected_coordinate}
    Let $R$ be a ring and $i\in I$.
    Then $\apply{\freeBasisVector{R}{I}{i}}{i}=\ringOne[R]$.
\end{proposition}
\begin{proof}
    Follows by
        \auto{free_basis_vector_point_mass,
            ring_zero_type,ring_one_type,point_mass_selected}.
\end{proof}

\begin{proposition}\label{free_basis_vector_other_coordinate}
    Let $R$ be a ring and $i,j\in I$.
    Suppose $j\neq i$.
    Then $\apply{\freeBasisVector{R}{I}{i}}{j}=\ringZero[R]$.
\end{proposition}
\begin{proof}
    Follows by
        \auto{free_basis_vector_point_mass,
            ring_zero_type,ring_one_type,point_mass_other}.
\end{proof}

\begin{proposition}\label{free_basis_vector_support}
    Let $R$ be a ring and $i\in I$.
    Then
        $(\ringOne[R]=\ringZero[R]\implies
            \supportOn{I}{\freeBasisVector{R}{I}{i}}{\ringZero[R]}
                =\emptyset)$ and
        $(\ringOne[R]\neq\ringZero[R]\implies
            \supportOn{I}{\freeBasisVector{R}{I}{i}}{\ringZero[R]}
                =\{i\})$.
\end{proposition}
\begin{proof}
    Let $p=\pointMass{I}{\carrier[R]}{\ringZero[R]}{i}{
        \ringOne[R]}$.
    Let $C=I\times\{\carrier[R]\}$ and
        $z=I\times\{\ringZero[R]\}$.
    $\ringZero[R],\ringOne[R]\in\carrier[R]$ by
        \obv{ring_zero_type,ring_one_type}.
    $p\in\funs{I}{\carrier[R]}$ by
        \auto{point_mass_finitely_supported,
            finitely_supported_functions_iff}.
    $\dependentSupport{I}{p}{z}=\supportOn{I}{p}{\ringZero[R]}$ by
        \auto{constant_family_support}.
    $C$ is a function on $I$ by
        \auto{constant_graph_type,funs_is_function,funs}.
    $z\in\dependentProduct{I}{C}$ by
        \auto{constant_family_dependent_product,
            constant_graph_type,singleton_subset_intro,
            funs_weaken_codom}.
    $C(i)=\carrier[R]$ and $z(i)=\ringZero[R]$ by
        \auto{constant_graph_application}.
    Follows by
        \auto{free_basis_vector_point_mass,point_mass,
            dependent_point_mass_support}.
\end{proof}

\begin{proposition}\label{free_module_point_mass_contribution_apply}
    Let $R$ be a ring.
    Let $F=I\times\{\LeftRegularModule{R}\}$ and
        $A=\moduleAdditiveGroupFamily{I}{F}$.
    Let $S=\FreeModule{R}{I}$ and
        $H=\ModuleAdditiveGroup{S}$.
    Let $x\in\carrier[S]$.
    Let $u=\directSumGroupInclusionFamily{I}{A}$.
    Let $i\in\supportOn{I}{x}{\ringZero[R]}$.
    Then
        $\apply{\directSumGroupContribution{I}{A}{H}{u}{x}}{i}
            =\pointMass{I}{\carrier[R]}{\ringZero[R]}{i}{x(i)}$.
\end{proposition}
\begin{proof}
    Let $z=\productMonoidNeutral{I}{A}$.
    $F$ is a function on $I$ by
        \auto{free_module_regular_family_type,
            funs_is_function,funs}.
    For every $j\in I$ we have $F(j)$ is a left module over $R$ by
        \auto{free_module_regular_family_apply,
            left_regular_module_is_exact}.
    Hence $A$ is a function on $I$ by
        \auto{module_additive_group_family_type}.
    For every $j\in I$ we have $A(j)$ is a commutative group by
        \auto{module_additive_group_family_apply,
            module_additive_group_is_exact_commutative_group}.
    Hence for every $j\in I$ we have $A(j)$ is a group by
        \obv{commutative_group_is_group}.
    $H=\DirectSumGroup{I}{A}$ by
        \auto{free_module_constructor_eq,
            direct_sum_module_additive_group_exact}.
    Hence $H$ is a commutative group by
        \auto{direct_sum_group_is_exact}.
    Hence $H$ is a group by \obv{commutative_group_is_group}.
    $u$ is a function on $I$ by
        \auto{direct_sum_group_inclusion_family_type}.
    For every $j\in I$ we have $\groupHom{u(j)}{A(j)}{H}$ by
        \auto{direct_sum_group_inclusion_family_apply,
            direct_sum_group_inclusion_hom}.
    $x\in\carrier[\DirectSumGroup{I}{A}]$ by
        \auto{free_module_constructor_eq,direct_sum_module_carrier}.
    $z=I\times\{\ringZero[R]\}$ by
        \obv{free_module_zero_section}.
    $x\in\finitelySupportedFunctions
        {I}{\carrier[R]}{\ringZero[R]}$ by
        \auto{free_module_carrier}.
    Hence $x\in\funs{I}{\carrier[R]}$ by
        \obv{finitely_supported_functions}.
    $\dependentSupport{I}{x}{I\times\{\ringZero[R]\}}
            =\supportOn{I}{x}{\ringZero[R]}$ by
        \auto{ring_zero_type,constant_family_support}.
    Hence $\dependentSupport{I}{x}{z}
        =\supportOn{I}{x}{\ringZero[R]}$.
    Hence $i\in\dependentSupport{I}{x}{z}$.
    Hence $i\in I$ by \obv{dependent_support}.
    $\apply{\directSumGroupContribution{I}{A}{H}{u}{x}}{i}
        =\apply{u(i)}{x(i)}$ by
        \auto{direct_sum_group_contribution_apply}.
    $u(i)=\directSumGroupInclusion{I}{A}{i}$ by
        \auto{direct_sum_group_inclusion_family_apply}.
    $\directSumGroupInclusion{I}{A}{i}
        =\directSumModuleInclusion{R}{I}{F}{i}$ by
        \obv{direct_sum_module_inclusion}.
    $x(i)\in\carrier[F(i)]$ by
        \auto{direct_sum_group_carrier,
            dependent_finitely_supported_functions_iff,
            dependent_product_coordinate,
            product_monoid_carrier_family_apply,
            module_additive_group_family_apply,
            module_additive_group_fields}.
    Hence
        $\apply{\directSumModuleInclusion{R}{I}{F}{i}}{x(i)}
            =\dependentPointMass{I}{
                \productMonoidCarrierFamily{I}{A}}{z}{i}{x(i)}$ by
        \obv{direct_sum_module_inclusion_apply}.
    $\productMonoidCarrierFamily{I}{A}
        =I\times\{\carrier[R]\}$ by
        \obv{free_module_carrier_family}.
    Follows by \auto{point_mass}.
\end{proof}

\begin{theorem}\label{free_module_point_mass_decomposition}
    Let $R$ be a ring.
    Let $F=I\times\{\LeftRegularModule{R}\}$ and
        $A=\moduleAdditiveGroupFamily{I}{F}$.
    Let $S=\FreeModule{R}{I}$ and
        $H=\ModuleAdditiveGroup{S}$.
    Let $x\in\carrier[S]$ and
        $u=\directSumGroupInclusionFamily{I}{A}$.
    Then
        $\finiteFold{H}{\supportOn{I}{x}{\ringZero[R]}}{
            \directSumGroupContribution{I}{A}{H}{u}{x}}=x$.
\end{theorem}
\begin{proof}
    Let $z=\productMonoidNeutral{I}{A}$.
    $F$ is a function on $I$ by
        \auto{free_module_regular_family_type,
            funs_is_function,funs}.
    For every $i\in I$ we have $F(i)$ is a left module over $R$ by
        \auto{free_module_regular_family_apply,
            left_regular_module_is_exact}.
    Hence $A$ is a function on $I$ by
        \auto{module_additive_group_family_type}.
    For every $i\in I$ we have $A(i)$ is a commutative group by
        \auto{module_additive_group_family_apply,
            module_additive_group_is_exact_commutative_group}.
    $H=\DirectSumGroup{I}{A}$ by
        \auto{free_module_constructor_eq,
            direct_sum_module_additive_group_exact}.
    $x\in\carrier[\DirectSumGroup{I}{A}]$ by
        \auto{free_module_constructor_eq,direct_sum_module_carrier}.
    $z=I\times\{\ringZero[R]\}$ by
        \obv{free_module_zero_section}.
    $x\in\finitelySupportedFunctions
        {I}{\carrier[R]}{\ringZero[R]}$ by
        \auto{free_module_carrier}.
    Hence $x\in\funs{I}{\carrier[R]}$ by
        \obv{finitely_supported_functions}.
    $\dependentSupport{I}{x}{I\times\{\ringZero[R]\}}
            =\supportOn{I}{x}{\ringZero[R]}$ by
        \auto{ring_zero_type,constant_family_support}.
    Hence $\dependentSupport{I}{x}{z}
        =\supportOn{I}{x}{\ringZero[R]}$.
    $\finiteFold{\DirectSumGroup{I}{A}}{
        \dependentSupport{I}{x}{z}}{
        \directSumGroupContribution{I}{A}{
            \DirectSumGroup{I}{A}}{u}{x}}=x$ by
        \obv{direct_sum_group_point_mass_decomposition}.
\end{proof}
\subsection{Empty, singleton, and reindexed free modules}

\begin{theorem}\label{empty_free_module_eq_zero_module}
    Let $R$ be a ring.
    Then $\FreeModule{R}{\emptyset}=\ZeroModule{R}$.
\end{theorem}
\begin{proof}
    Let $F=\FreeModule{R}{\emptyset}$ and $Z=\ZeroModule{R}$.
    $F$ is an exact left module over $R$ by
        \auto{free_module_is_exact}.
    $Z$ is an exact left module over $R$ by
        \auto{zero_module_is_exact}.
    $\carrier[F]
        =\finitelySupportedFunctions{\emptyset}{\carrier[R]}{
            \ringZero[R]}$ by \auto{free_module_carrier}.
    Show that
        $\finitelySupportedFunctions{\emptyset}{\carrier[R]}{
            \ringZero[R]}=\{\emptyset\}$.
    \begin{subproof}
        Show that every element of
            $\finitelySupportedFunctions{\emptyset}{\carrier[R]}{
                \ringZero[R]}$ is an element of $\{\emptyset\}$.
        \begin{subproof}
            Fix $x\in\finitelySupportedFunctions{\emptyset}{
                \carrier[R]}{\ringZero[R]}$.
            Hence $x\in\funs{\emptyset}{\carrier[R]}$ by
                \obv{finitely_supported_functions}.
            Hence $x\in\{\emptyset\}$ by
                \auto{empty_domain_function_space}.
        \end{subproof}
        Show that every element of $\{\emptyset\}$ is an element of
            $\finitelySupportedFunctions{\emptyset}{\carrier[R]}{
                \ringZero[R]}$.
        \begin{subproof}
            Fix $x\in\{\emptyset\}$.
            Hence $x=\emptyset$ by \auto{singleton_iff}.
            $\emptyset\in\funs{\emptyset}{\carrier[R]}$ by
                \auto{empty_domain_function_space,singleton_iff}.
            $\supportOn{\emptyset}{\emptyset}{\ringZero[R]}
                =\emptyset$ by \auto{setext,support_on_iff,emptyset}.
            Hence $\supportOn{\emptyset}{\emptyset}{\ringZero[R]}$
                is finite by \obv{emptyset_is_finite}.
            Follows by
                \auto{finitely_supported_functions_iff}.
        \end{subproof}
        Follows by \obv{setext}.
    \end{subproof}
    Hence $\carrier[F]=\{\emptyset\}$.
    $\carrier[Z]=\{\emptyset\}$ by \auto{zero_module_fields}.
    Hence $\carrier[F]=\carrier[Z]$.
    $\ModuleAdditiveGroup{F}$ is a group by
        \obv{module_additive_group_is_group}.
    Hence $\mul[F]\in
        \funs{\{\emptyset\}\times\{\emptyset\}}{\{\emptyset\}}$ by
        \auto{module_additive_group_fields,
            group_is_monoid,monoid_is_magma,magma_mul_graph}.
    $\funs{\{\emptyset\}\times\{\emptyset\}}{\{\emptyset\}}
        =\{(\{\emptyset\}\times\{\emptyset\})
            \times\{\emptyset\}\}$ by
        \obv{singleton_codom_function_space}.
    Hence $\mul[F]=(\{\emptyset\}\times\{\emptyset\})
        \times\{\emptyset\}$ by \auto{singleton_iff}.
    $\mul[Z]=(\{\emptyset\}\times\{\emptyset\})
        \times\{\emptyset\}$ by \auto{zero_module_fields}.
    Hence $\mul[F]=\mul[Z]$.
    $\neutral[F]\in\carrier[F]$ by
        \auto{module_additive_group_fields,subgroup_neutral,
            full_carrier_subgroup}.
    Hence $\neutral[F]=\emptyset$ by \auto{singleton_iff}.
    $\neutral[Z]=\emptyset$ by \auto{zero_module_fields}.
    Hence $\neutral[F]=\neutral[Z]$.
    $\groupInv[F]\in\funs{\{\emptyset\}}{\{\emptyset\}}$ by
        \auto{module_additive_group_fields,group_inverse_graph}.
    $\funs{\{\emptyset\}}{\{\emptyset\}}
        =\{\{\emptyset\}\times\{\emptyset\}\}$ by
        \obv{singleton_codom_function_space}.
    Hence $\groupInv[F]=\{\emptyset\}\times\{\emptyset\}$ by
        \auto{singleton_iff}.
    $\groupInv[Z]=\{\emptyset\}\times\{\emptyset\}$ by
        \auto{zero_module_fields}.
    Hence $\groupInv[F]=\groupInv[Z]$.
    $\scalarMul[F]\in
        \funs{\carrier[R]\times\{\emptyset\}}{\{\emptyset\}}$ by
        \auto{module_scalar_graph}.
    $\funs{\carrier[R]\times\{\emptyset\}}{\{\emptyset\}}
        =\{(\carrier[R]\times\{\emptyset\})
            \times\{\emptyset\}\}$ by
        \obv{singleton_codom_function_space}.
    Hence $\scalarMul[F]=(\carrier[R]\times\{\emptyset\})
        \times\{\emptyset\}$ by \auto{singleton_iff}.
    $\scalarMul[Z]=(\carrier[R]\times\{\emptyset\})
        \times\{\emptyset\}$ by \obv{zero_module_fields}.
    Hence $\scalarMul[F]=\scalarMul[Z]$.
    Follows by \auto{exact_left_module_extensionality}.
\end{proof}

\begin{theorem}\label{singleton_free_module_inclusion_isomorphism}
    Let $R$ be a ring.
    Let $F=\{i\}\times\{\LeftRegularModule{R}\}$ and
        $u=\directSumModuleInclusion{R}{\{i\}}{F}{i}$.
    Then $u$ is a module isomorphism from $\LeftRegularModule{R}$ to
        $\FreeModule{R}{\{i\}}$ over $R$.
\end{theorem}
\begin{proof}
    Let $L=\LeftRegularModule{R}$ and
        $S=\FreeModule{R}{\{i\}}$.
    $F$ is a function on $\{i\}$ by
        \auto{free_module_regular_family_type,
            funs_is_function,funs}.
    For every $j\in\{i\}$ we have $F(j)$ is a left module over $R$ by
        \auto{free_module_regular_family_apply,
            left_regular_module_is_exact}.
    $F(i)=L$ by
        \auto{free_module_regular_family_apply,singleton_iff}.
    $S=\DirectSumModule{R}{\{i\}}{F}$ by
        \obv{free_module_constructor_eq}.
    Hence $u$ is a linear map from $L$ to $S$ over $R$ by
        \auto{direct_sum_module_inclusion_linear}.
    $u\in\funs{\carrier[L]}{\carrier[S]}$ by
        \obv{linear_map_type}.
    Hence $u$ is a function by \obv{funs_is_function}.
    Show that $u$ is injective.
    \begin{subproof}
        Show that for all $a,b\in\dom{u}$ such that $u(a)=u(b)$
            we have $a=b$.
        \begin{subproof}
            Fix $a,b\in\dom{u}$.
            Assume $u(a)=u(b)$.
            Hence $a,b\in\carrier[L]$ by
                \auto{linear_map_type,funs}.
            Let $A=\moduleAdditiveGroupFamily{\{i\}}{F}$ and
                $C=\productMonoidCarrierFamily{\{i\}}{A}$.
            Let $z=\productMonoidNeutral{\{i\}}{A}$.
            $u(a)=\dependentPointMass{\{i\}}{C}{z}{i}{a}$ and
                $u(b)=\dependentPointMass{\{i\}}{C}{z}{i}{b}$ by
                \auto{direct_sum_module_inclusion_apply}.
            $A$ is a function on $\{i\}$ by
                \auto{module_additive_group_family_type}.
            $C$ is a function on $\{i\}$ by
                \auto{product_monoid_carrier_family_type}.
            For every $j\in\{i\}$ we have $A(j)$ is a commutative group by
                \auto{module_additive_group_family_apply,
                    module_additive_group_is_exact_commutative_group}.
            Hence for every $j\in\{i\}$ we have $A(j)$ is a monoid by
                \obv{commutative_group_is_group,group_is_monoid}.
            $z\in\dependentProduct{\{i\}}{C}$ by
                \auto{product_monoid_neutral_type}.
            $a,b\in C(i)$ by
                \auto{product_monoid_carrier_family_apply,
                    module_additive_group_family_apply,
                    module_additive_group_fields}.
            Hence $(u(a))(i)=a$ and $(u(b))(i)=b$ by
                \auto{dependent_point_mass_selected}.
        \end{subproof}
        Follows by \auto{injective_function}.
    \end{subproof}
    Show that $u\in\surj{\carrier[L]}{\carrier[S]}$.
    \begin{subproof}
        Show that for every $x\in\carrier[S]$ there exists
            $a\in\carrier[L]$ such that $u(a)=x$.
        \begin{subproof}
            Fix $x\in\carrier[S]$.
            Let $A=\moduleAdditiveGroupFamily{\{i\}}{F}$ and
                $C=\productMonoidCarrierFamily{\{i\}}{A}$.
            Let $z=\productMonoidNeutral{\{i\}}{A}$.
            $x\in\carrier[\DirectSumGroup{\{i\}}{A}]$ by
                \auto{free_module_constructor_eq,
                    direct_sum_module_carrier}.
            Hence $x\in\dependentProduct{\{i\}}{C}$ by
                \auto{direct_sum_group_carrier,
                    dependent_finitely_supported_functions_iff}.
            Hence $x(i)\in C(i)$ by
                \auto{dependent_product_coordinate,singleton_iff}.
            Hence $x(i)\in\carrier[L]$ by
                \auto{product_monoid_carrier_family_apply,
                    module_additive_group_family_apply,
                    module_additive_group_fields}.
            $u(x(i))\in\carrier[S]$ by
                \obv{linear_map_apply_type}.
            $u(x(i))=\dependentPointMass{\{i\}}{C}{z}{i}{x(i)}$ by
                \auto{direct_sum_module_inclusion_apply}.
            Hence $u(x(i))\in\dependentProduct{\{i\}}{C}$ by
                \auto{direct_sum_group_carrier,
                    dependent_finitely_supported_functions_iff,
                    free_module_constructor_eq,
                    direct_sum_module_carrier}.
            Show that for every $j\in\{i\}$ we have
                $(u(x(i)))(j)=x(j)$.
            \begin{subproof}
                Fix $j\in\{i\}$.
                Hence $j=i$ by \auto{singleton_iff}.
                $C$ is a function on $\{i\}$ by
                    \auto{module_additive_group_family_type,
                        product_monoid_carrier_family_type}.
                $A$ is a function on $\{i\}$ by
                    \auto{module_additive_group_family_type}.
                For every $k\in\{i\}$ we have $A(k)$ is a commutative group by
                    \auto{module_additive_group_family_apply,
                        module_additive_group_is_exact_commutative_group}.
                Hence for every $k\in\{i\}$ we have $A(k)$ is a monoid by
                    \obv{commutative_group_is_group,group_is_monoid}.
                $z\in\dependentProduct{\{i\}}{C}$ by
                    \auto{product_monoid_neutral_type}.
                $(u(x(i)))(i)=x(i)$ by
                    \auto{dependent_point_mass_selected}.
            \end{subproof}
            $u(x(i)),x\in\finitelySupportedFunctions{\{i\}}{
                \carrier[R]}{\ringZero[R]}$ by
                \auto{free_module_carrier}.
            Hence $u(x(i)),x\in\funs{\{i\}}{\carrier[R]}$ by
                \obv{finitely_supported_functions}.
            Hence $u(x(i))=x$ by
                \obv{typed_functions_equal_from_agreement}.
            Follows by \auto.
        \end{subproof}
        Follows by \obv{surj}.
    \end{subproof}
    Hence $u\in\bijections{\carrier[L]}{\carrier[S]}$ by
        \obv{bijections}.
    Follows by \auto{module_isomorphism}.
\end{proof}

\begin{proposition}\label{free_module_regular_family_reindex}
    Let $e\in\funs{J}{I}$.
    Then
        $(I\times\{\LeftRegularModule{R}\})\circ e
            =J\times\{\LeftRegularModule{R}\}$.
\end{proposition}
\begin{proof}
    Let $F=I\times\{\LeftRegularModule{R}\}$ and
        $K=J\times\{\LeftRegularModule{R}\}$.
    $F\in\funs{I}{\{\LeftRegularModule{R}\}}$ and
        $K\in\funs{J}{\{\LeftRegularModule{R}\}}$ by
        \auto{constant_graph_type}.
    Hence $F$ and $e$ are functions by
        \obv{funs_is_function}.
    $F$ is composable with $e$ by \auto{funs_ran,funs}.
    Hence $F\circ e\in\funs{J}{\{\LeftRegularModule{R}\}}$ by
        \obv{funs_circ}.
    Hence $F\circ e$ and $K$ are functions and
        $\dom{F\circ e}=J=\dom{K}$ by
        \auto{funs_is_function,funs}.
    Show that for every $j\in J$ we have $(F\circ e)(j)=K(j)$.
    \begin{subproof}
        Fix $j\in J$.
        $e$ is a function and $j\in\dom{e}$ by
            \auto{funs_is_function,funs}.
        $e(j)\in I$ by \obv{funs_type_apply}.
        $(F\circ e)(j)=F(e(j))$ by \obv{circ_apply}.
        $F(e(j))=\LeftRegularModule{R}$ and
            $K(j)=\LeftRegularModule{R}$ by
            \auto{constant_graph_application}.
    \end{subproof}
    Follows by \auto{funext_on_domain}.
\end{proof}

\begin{theorem}\label{bijective_free_module_reindex_isomorphism}
    Let $R$ be a ring.
    Let $e\in\bijections{J}{I}$.
    Let $F=I\times\{\LeftRegularModule{R}\}$ and
        $A=\moduleAdditiveGroupFamily{I}{F}$.
    Then $\directSumGroupReindex{I}{J}{A}{e}$ is a module isomorphism
        from $\FreeModule{R}{I}$ to $\FreeModule{R}{J}$ over $R$.
\end{theorem}
\begin{proof}
    $F$ is a function on $I$ by
        \auto{free_module_regular_family_type,
            funs_is_function,funs}.
    For every $i\in I$ we have $F(i)$ is a left module over $R$ by
        \auto{free_module_regular_family_apply,
            left_regular_module_is_exact}.
    $\directSumGroupReindex{I}{J}{A}{e}$ is a module isomorphism from
        $\DirectSumModule{R}{I}{F}$ to
        $\DirectSumModule{R}{J}{F\circ e}$ over $R$ by
        \obv{bijective_direct_sum_module_reindex_isomorphism}.
    $e\in\funs{J}{I}$ by \obv{bijections_to_funs}.
    $F\circ e=J\times\{\LeftRegularModule{R}\}$ by
        \auto{free_module_regular_family_reindex}.
    Follows by \auto{free_module_constructor_eq}.
\end{proof}

\begin{proposition}\label{free_module_reindex_apply}
    Let $R$ be a ring.
    Let $e\in\bijections{J}{I}$.
    Let $F=I\times\{\LeftRegularModule{R}\}$ and
        $A=\moduleAdditiveGroupFamily{I}{F}$.
    Let $x\in\carrier[\FreeModule{R}{I}]$.
    Then
        $\apply{\directSumGroupReindex{I}{J}{A}{e}}{x}=x\circ e$.
\end{proposition}
\begin{proof}
    $F$ is a function on $I$ by
        \auto{free_module_regular_family_type,
            funs_is_function,funs}.
    For every $i\in I$ we have $F(i)$ is a left module over $R$ by
        \auto{free_module_regular_family_apply,
            left_regular_module_is_exact}.
    Hence $A$ is a function on $I$ by
        \auto{module_additive_group_family_type}.
    For every $i\in I$ we have $A(i)$ is a commutative group by
        \auto{module_additive_group_family_apply,
            module_additive_group_is_exact_commutative_group}.
    $e\in\funs{J}{I}$ by \obv{bijections_to_funs}.
    For every $i\in I$ we have $\preimg{e}{\{i\}}$ is finite by
        \obv{bijection_has_finite_fibers}.
    $x\in\carrier[\DirectSumGroup{I}{A}]$ by
        \auto{free_module_constructor_eq,direct_sum_module_carrier}.
    Follows by \obv{direct_sum_group_reindex_apply}.
\end{proof}

\begin{proposition}\label{free_basis_vector_reindex}
    Let $R$ be a ring.
    Let $e\in\bijections{J}{I}$ and $j\in J$.
    Let $F=I\times\{\LeftRegularModule{R}\}$ and
        $A=\moduleAdditiveGroupFamily{I}{F}$.
    Let $r=\directSumGroupReindex{I}{J}{A}{e}$.
    Then $r(\freeBasisVector{R}{I}{e(j)})
        =\freeBasisVector{R}{J}{j}$.
\end{proposition}
\begin{proof}
    Let $x=\freeBasisVector{R}{I}{e(j)}$ and
        $y=\freeBasisVector{R}{J}{j}$.
    $e\in\funs{J}{I}$ by \obv{bijections_to_funs}.
    Hence $e$ is a function and $j\in\dom{e}$ by
        \auto{funs_is_function,funs}.
    Hence $e(j)\in I$ by \obv{funs_type_apply}.
    $x\in\carrier[\FreeModule{R}{I}]$ and
        $y\in\carrier[\FreeModule{R}{J}]$ by
        \auto{free_basis_vector_type}.
    $r(x)=x\circ e$ by \auto{free_module_reindex_apply}.
    Hence $r(x)\in\carrier[\FreeModule{R}{J}]$ by
        \auto{bijective_free_module_reindex_isomorphism,
            module_isomorphism_is_linear_map,linear_map_apply_type}.
    Hence $r(x),y\in\funs{J}{\carrier[R]}$ by
        \auto{free_module_carrier,
            finitely_supported_functions_iff}.
    Show that for every $k\in J$ we have $(r(x))(k)=y(k)$.
    \begin{subproof}
        Fix $k\in J$.
        $e$ is a function and $k\in\dom{e}$ by
            \auto{funs_is_function,funs}.
        $x$ is a function by
            \auto{free_module_carrier,
                finitely_supported_functions_iff,funs_is_function}.
        $x$ is composable with $e$ by
            \auto{free_module_carrier,
                finitely_supported_functions_iff,funs_ran,funs}.
        $(r(x))(k)=(x\circ e)(k)$.
        Hence $(r(x))(k)=x(e(k))$ by \auto{circ_apply}.
        \begin{byCase}
        \caseOf{$k=j$.}
            Hence $e(k)=e(j)$.
            $x(e(j))=\ringOne[R]$ and $y(j)=\ringOne[R]$ by
                \auto{free_basis_vector_selected_coordinate}.
        \caseOf{$k\neq j$.}
            $e$ is injective by \obv{bijections}.
            $k,j\in\dom{e}$ by \obv{funs}.
            $e(k)\neq e(j)$ by
                \auto{injective_function}.
            $e(k)\in I$ by \obv{funs_type_apply}.
            $x(e(k))=\ringZero[R]$ by
                \auto{free_basis_vector_other_coordinate}.
            $y(k)=\ringZero[R]$ by
                \auto{free_basis_vector_other_coordinate}.
        \end{byCase}
    \end{subproof}
    Follows by \auto{typed_functions_equal_from_agreement}.
\end{proof}

\subsection{The universal property}

For a vector $x$ in a left module, scalar multiplication gives the unique
linear map from the left regular module that sends its unit to $x$.  These are
the component maps assembled by the free-module lift.

\begin{theorem}\label{left_regular_scalar_extension_linear}
    Let $R$ be a ring.
    Let $M$ be a left module over $R$ and $x\in\carrier[M]$.
    Let
        $f=\Apply{\graphOnOp}{
            \carrier[R],\carrier[M],
            \fun{a:\Set}{\scalarMul[M](a,x)}}$.
    Then $f$ is a linear map from $\LeftRegularModule{R}$ to $M$ over $R$.
\end{theorem}
\begin{proof}
    Let $L=\LeftRegularModule{R}$.
    $L$ is a left module over $R$ by
        \auto{left_regular_module_is_exact}.
    For every $a\in\carrier[R]$ we have
        $\scalarMul[M](a,x)\in\carrier[M]$ by
        \obv{module_scalar_type}.
    Hence $f\in\funs{\carrier[R]}{\carrier[M]}$ by
        \auto{logical_graph_on_function}.
    $\carrier[L]=\carrier[R]$ and
        $\mul[L]=\ringAdd[R]$ and
        $\scalarMul[L]=\ringMul[R]$ by
        \auto{left_regular_module_fields,additive_group_fields}.
    Hence $f\in\funs{\carrier[L]}{\carrier[M]}$.
    Show that for every $a\in\carrier[R]$ we have
        $f(a)=\scalarMul[M](a,x)$.
    \begin{subproof}
        Follows by \auto{logical_graph_on_apply}.
    \end{subproof}
    Show that for all $a,b\in\carrier[L]$ we have
        $f(\mul[L](a,b))=\mul[M](f(a),f(b))$.
    \begin{subproof}
        Fix $a,b\in\carrier[L]$.
        Hence $a,b\in\carrier[R]$.
        $\ringAdd[R](a,b)\in\carrier[R]$ by \obv{ring_add_type}.
        $f(\mul[L](a,b))
            =\scalarMul[M](\ringAdd[R](a,b),x)$ by \auto.
        $f(a)=\scalarMul[M](a,x)$ and
            $f(b)=\scalarMul[M](b,x)$.
        Follows by \auto{module_scalar_add_distributive}.
    \end{subproof}
    Show that for all $a,b$ such that $a\in\carrier[R]$ and
        $b\in\carrier[L]$ we have
        $f(\scalarMul[L](a,b))=\scalarMul[M](a,f(b))$.
    \begin{subproof}
        Fix $a,b$.
        Assume $a\in\carrier[R]$ and $b\in\carrier[L]$.
        Hence $b\in\carrier[R]$.
        $\ringMul[R](a,b)\in\carrier[R]$ by \obv{ring_mul_type}.
        $f(\scalarMul[L](a,b))
            =\scalarMul[M](\ringMul[R](a,b),x)$ by \auto.
        $f(b)=\scalarMul[M](b,x)$.
        Follows by \auto{module_scalar_associative}.
    \end{subproof}
    Follows by \auto{linear_map_expanded_iff}.
\end{proof}

\begin{proposition}\label{free_module_coefficient_family_type}
    Let $R$ be a ring.
    Let $M$ be a left module over $R$.
    Let $h\in\funs{I}{\carrier[M]}$.
    Let
        $g=\Apply{\graphOnDomainOp}{
            I,\fun{j:\Set}{
                \Apply{\graphOnOp}{
                    \carrier[R],\carrier[M],
                    \fun{a:\Set}{\scalarMul[M](a,h(j))}}}}$.
    Then $g$ is a function on $I$ and
        for every $i\in I$ we have $g(i)$ is a linear map from
            $\LeftRegularModule{R}$ to $M$ over $R$.
\end{proposition}
\begin{proof}
    $g\in\funs{I}{\ran{g}}$ by
        \auto{logical_graph_on_domain_function}.
    Hence $g$ is a function on $I$ by
        \obv{funs_is_function,funs}.
    Show that for every $i\in I$ we have $g(i)$ is a linear map from
        $\LeftRegularModule{R}$ to $M$ over $R$.
    \begin{subproof}
        Fix $i\in I$.
        $h(i)\in\carrier[M]$ by \obv{funs_type_apply}.
        $g(i)=\Apply{\graphOnOp}{
            \carrier[R],\carrier[M],
            \fun{a:\Set}{\scalarMul[M](a,h(i))}}$ by
            \auto{logical_graph_on_domain_apply}.
        Follows by \obv{left_regular_scalar_extension_linear}.
    \end{subproof}
\end{proof}

\begin{proposition}\label{free_module_coefficient_family_apply}
    Let $R$ be a ring.
    Let $M$ be a left module over $R$.
    Let $h\in\funs{I}{\carrier[M]}$.
    Let
        $g=\Apply{\graphOnDomainOp}{
            I,\fun{j:\Set}{
                \Apply{\graphOnOp}{
                    \carrier[R],\carrier[M],
                    \fun{b:\Set}{\scalarMul[M](b,h(j))}}}}$.
    Let $i\in I$ and $a\in\carrier[R]$.
    Then $\apply{g(i)}{a}=\scalarMul[M](a,h(i))$.
\end{proposition}
\begin{proof}
    $g(i)=\Apply{\graphOnOp}{
        \carrier[R],\carrier[M],
        \fun{b:\Set}{\scalarMul[M](b,h(i))}}$ by
        \auto{logical_graph_on_domain_apply}.
    $h(i)\in\carrier[M]$ by \obv{funs_type_apply}.
    For every $b\in\carrier[R]$ we have
        $\scalarMul[M](b,h(i))\in\carrier[M]$ by
        \obv{module_scalar_type}.
    Follows by \auto{logical_graph_on_apply}.
\end{proof}

\begin{definition}\label{free_module_lift}
    $\freeModuleLift{R}{I}{M}{h}
        =\directSumModuleMediator{R}{I}{
            I\times\{\LeftRegularModule{R}\}}{M}{
            \Apply{\graphOnDomainOp}{
                I,\fun{i:\Set}{
                    \Apply{\graphOnOp}{
                        \carrier[R],\carrier[M],
                        \fun{a:\Set}{\scalarMul[M](a,h(i))}}}}}$.
\end{definition}

\begin{proposition}\label{free_module_lift_type}
    Let $R$ be a ring.
    Let $M$ be a left module over $R$.
    Let $h\in\funs{I}{\carrier[M]}$.
    Then
        $\freeModuleLift{R}{I}{M}{h}\in
            \funs{\carrier[\FreeModule{R}{I}]}{\carrier[M]}$.
\end{proposition}
\begin{proof}
    Let $F=I\times\{\LeftRegularModule{R}\}$.
    Let
        $g=\Apply{\graphOnDomainOp}{
            I,\fun{i:\Set}{
                \Apply{\graphOnOp}{
                    \carrier[R],\carrier[M],
                    \fun{a:\Set}{\scalarMul[M](a,h(i))}}}}$.
    $F$ is a function on $I$ by
        \auto{free_module_regular_family_type,
            funs_is_function,funs}.
    For every $i\in I$ we have $F(i)$ is a left module over $R$ by
        \auto{free_module_regular_family_apply,
            left_regular_module_is_exact}.
    $g$ is a function on $I$ and for every $i\in I$ we have
        $g(i)$ is a linear map from $\LeftRegularModule{R}$ to $M$ over $R$ by
        \obv{free_module_coefficient_family_type}.
    Hence for every $i\in I$ we have $g(i)$ is a linear map from
        $F(i)$ to $M$ over $R$ by
        \auto{free_module_regular_family_apply}.
    $\directSumModuleMediator{R}{I}{F}{M}{g}\in
        \funs{\carrier[\DirectSumModule{R}{I}{F}]}{\carrier[M]}$ by
        \obv{direct_sum_module_mediator_type}.
    Follows by
        \auto{free_module_lift,free_module_constructor_eq}.
\end{proof}

\begin{theorem}\label{free_module_lift_linear}
    Let $R$ be a ring.
    Let $M$ be a left module over $R$.
    Let $h\in\funs{I}{\carrier[M]}$.
    Then $\freeModuleLift{R}{I}{M}{h}$ is a linear map from
        $\FreeModule{R}{I}$ to $M$ over $R$.
\end{theorem}
\begin{proof}
    Let $F=I\times\{\LeftRegularModule{R}\}$.
    Let
        $g=\Apply{\graphOnDomainOp}{
            I,\fun{i:\Set}{
                \Apply{\graphOnOp}{
                    \carrier[R],\carrier[M],
                    \fun{a:\Set}{\scalarMul[M](a,h(i))}}}}$.
    $F$ is a function on $I$ by
        \auto{free_module_regular_family_type,
            funs_is_function,funs}.
    For every $i\in I$ we have $F(i)$ is a left module over $R$ by
        \auto{free_module_regular_family_apply,
            left_regular_module_is_exact}.
    $g$ is a function on $I$ and for every $i\in I$ we have
        $g(i)$ is a linear map from $\LeftRegularModule{R}$ to $M$ over $R$ by
        \obv{free_module_coefficient_family_type}.
    Hence for every $i\in I$ we have $g(i)$ is a linear map from
        $F(i)$ to $M$ over $R$ by
        \auto{free_module_regular_family_apply}.
    $\directSumModuleMediator{R}{I}{F}{M}{g}$ is a linear map from
        $\DirectSumModule{R}{I}{F}$ to $M$ over $R$ by
        \obv{direct_sum_module_mediator_linear}.
    Follows by
        \auto{free_module_lift,free_module_constructor_eq}.
\end{proof}

\begin{theorem}\label{free_module_lift_factorization}
    Let $R$ be a ring.
    Let $M$ be a left module over $R$.
    Let $h\in\funs{I}{\carrier[M]}$.
    Let
        $g=\Apply{\graphOnDomainOp}{
            I,\fun{j:\Set}{
                \Apply{\graphOnOp}{
                    \carrier[R],\carrier[M],
                    \fun{a:\Set}{\scalarMul[M](a,h(j))}}}}$.
    Let $i\in I$.
    Then
        $\freeModuleLift{R}{I}{M}{h}\circ
            \directSumModuleInclusion{R}{I}{
                I\times\{\LeftRegularModule{R}\}}{i}=g(i)$.
\end{theorem}
\begin{proof}
    Let $F=I\times\{\LeftRegularModule{R}\}$.
    $F$ is a function on $I$ by
        \auto{free_module_regular_family_type,
            funs_is_function,funs}.
    For every $j\in I$ we have $F(j)$ is a left module over $R$ by
        \auto{free_module_regular_family_apply,
            left_regular_module_is_exact}.
    $g$ is a function on $I$ and for every $j\in I$ we have
        $g(j)$ is a linear map from $\LeftRegularModule{R}$ to $M$ over $R$ by
        \obv{free_module_coefficient_family_type}.
    Hence for every $j\in I$ we have $g(j)$ is a linear map from
        $F(j)$ to $M$ over $R$ by
        \auto{free_module_regular_family_apply}.
    $\directSumModuleMediator{R}{I}{F}{M}{g}\circ
        \directSumModuleInclusion{R}{I}{F}{i}=g(i)$ by
        \obv{direct_sum_module_mediator_factorization}.
    Follows by \auto{free_module_lift}.
\end{proof}

\begin{theorem}\label{free_module_lift_basis_apply}
    Let $R$ be a ring.
    Let $M$ be a left module over $R$.
    Let $h\in\funs{I}{\carrier[M]}$ and $i\in I$.
    Then
        $\apply{\freeModuleLift{R}{I}{M}{h}}{
            \freeBasisVector{R}{I}{i}}=h(i)$.
\end{theorem}
\begin{proof}
    Let $F=I\times\{\LeftRegularModule{R}\}$ and
        $u=\directSumModuleInclusion{R}{I}{F}{i}$.
    Let $m=\freeModuleLift{R}{I}{M}{h}$ and
        $L=\LeftRegularModule{R}$.
    Let $S=\FreeModule{R}{I}$.
    Let
        $g=\Apply{\graphOnDomainOp}{
            I,\fun{j:\Set}{
                \Apply{\graphOnOp}{
                    \carrier[R],\carrier[M],
                    \fun{a:\Set}{\scalarMul[M](a,h(j))}}}}$.
    $m\circ u=g(i)$ by
        \auto{free_module_lift_factorization}.
    $m$ is a linear map from $S$ to $M$ over $R$ by
        \auto{free_module_lift_linear}.
    $F$ is a function on $I$ by
        \auto{free_module_regular_family_type,
            funs_is_function,funs}.
    For every $j\in I$ we have $F(j)$ is a left module over $R$ by
        \auto{free_module_regular_family_apply,
            left_regular_module_is_exact}.
    $F(i)=L$ and $S=\DirectSumModule{R}{I}{F}$ by
        \auto{free_module_regular_family_apply,
            free_module_constructor_eq}.
    Hence $u$ is a linear map from $L$ to $S$ over $R$ by
        \auto{direct_sum_module_inclusion_linear}.
    Hence $m\in\funs{\carrier[S]}{\carrier[M]}$ and
        $u\in\funs{\carrier[L]}{\carrier[S]}$ by
        \obv{linear_map_type}.
    Hence $m$ and $u$ are functions by \auto{funs_is_function}.
    $m$ is composable with $u$ by \auto{funs_ran,funs}.
    $\ringOne[R]\in\dom{u}$ by
        \auto{funs,left_regular_module_carrier,ring_one_type}.
    Hence $(m\circ u)(\ringOne[R])=m(u(\ringOne[R]))$ by
        \obv{circ_apply}.
    $u(\ringOne[R])=\freeBasisVector{R}{I}{i}$ by
        \auto{free_basis_vector}.
    $\ringOne[R]\in\carrier[R]$ by \obv{ring_one_type}.
    $\apply{g(i)}{\ringOne[R]}
        =\scalarMul[M](\ringOne[R],h(i))$ by
        \obv{free_module_coefficient_family_apply}.
    $h(i)\in\carrier[M]$ by \obv{funs_type_apply}.
    Hence $\scalarMul[M](\ringOne[R],h(i))=h(i)$ by
        \obv{module_scalar_one}.
\end{proof}

\begin{proposition}\label{free_module_lift_contribution_apply}
    Let $R$ be a ring.
    Let $M$ be a left module over $R$.
    Let $h\in\funs{I}{\carrier[M]}$.
    Let $F=I\times\{\LeftRegularModule{R}\}$ and
        $A=\moduleAdditiveGroupFamily{I}{F}$.
    Let
        $g=\Apply{\graphOnDomainOp}{
            I,\fun{j:\Set}{
                \Apply{\graphOnOp}{
                    \carrier[R],\carrier[M],
                    \fun{a:\Set}{\scalarMul[M](a,h(j))}}}}$.
    Let $x\in\carrier[\FreeModule{R}{I}]$.
    Let $i\in\supportOn{I}{x}{\ringZero[R]}$.
    Then
        $\apply{\directSumGroupContribution{I}{A}{
            \ModuleAdditiveGroup{M}}{g}{x}}{i}
            =\scalarMul[M](x(i),h(i))$.
\end{proposition}
\begin{proof}
    Let $z=\productMonoidNeutral{I}{A}$ and
        $H=\ModuleAdditiveGroup{M}$.
    $F$ is a function on $I$ by
        \auto{free_module_regular_family_type,
            funs_is_function,funs}.
    For every $j\in I$ we have $F(j)$ is a left module over $R$ by
        \auto{free_module_regular_family_apply,
            left_regular_module_is_exact}.
    Hence $A$ is a function on $I$ by
        \auto{module_additive_group_family_type}.
    For every $j\in I$ we have $A(j)$ is a commutative group by
        \auto{module_additive_group_family_apply,
            module_additive_group_is_exact_commutative_group}.
    Hence for every $j\in I$ we have $A(j)$ is a group by
        \obv{commutative_group_is_group}.
    $H$ is a commutative group by
        \obv{module_additive_group_is_exact_commutative_group}.
    Hence $H$ is a group by \obv{commutative_group_is_group}.
    $g$ is a function on $I$ and for every $j\in I$ we have
        $g(j)$ is a linear map from $\LeftRegularModule{R}$ to $M$ over $R$ by
        \obv{free_module_coefficient_family_type}.
    Hence for every $j\in I$ we have $\groupHom{g(j)}{A(j)}{H}$ by
        \auto{free_module_regular_family_apply,
            module_additive_group_family_apply,
            linear_map_is_additive_group_hom}.
    $x\in\carrier[\DirectSumGroup{I}{A}]$ by
        \auto{free_module_constructor_eq,direct_sum_module_carrier}.
    $z=I\times\{\ringZero[R]\}$ by
        \obv{free_module_zero_section}.
    $x\in\funs{I}{\carrier[R]}$ by
        \auto{free_module_carrier,
            finitely_supported_functions_iff}.
    $\dependentSupport{I}{x}{I\times\{\ringZero[R]\}}
        =\supportOn{I}{x}{\ringZero[R]}$ by
        \auto{ring_zero_type,constant_family_support}.
    Hence $i\in\dependentSupport{I}{x}{z}$.
    Hence $i\in I$ by \obv{dependent_support}.
    $\apply{\directSumGroupContribution{I}{A}{H}{g}{x}}{i}
        =\apply{g(i)}{x(i)}$ by
        \auto{direct_sum_group_contribution_apply}.
    $x(i)\in\carrier[R]$ by \obv{funs_type_apply}.
    Follows by \auto{free_module_coefficient_family_apply}.
\end{proof}

\begin{theorem}\label{free_module_lift_apply}
    Let $R$ be a ring.
    Let $M$ be a left module over $R$.
    Let $h\in\funs{I}{\carrier[M]}$.
    Let $F=I\times\{\LeftRegularModule{R}\}$ and
        $A=\moduleAdditiveGroupFamily{I}{F}$.
    Let
        $g=\Apply{\graphOnDomainOp}{
            I,\fun{i:\Set}{
                \Apply{\graphOnOp}{
                    \carrier[R],\carrier[M],
                    \fun{a:\Set}{\scalarMul[M](a,h(i))}}}}$.
    Let $x\in\carrier[\FreeModule{R}{I}]$.
    Then
        $\apply{\freeModuleLift{R}{I}{M}{h}}{x}
            =\finiteFold{\ModuleAdditiveGroup{M}}{
                \supportOn{I}{x}{\ringZero[R]}}{
                \directSumGroupContribution{I}{A}{
                    \ModuleAdditiveGroup{M}}{g}{x}}$.
\end{theorem}
\begin{proof}
    Let $z=\productMonoidNeutral{I}{A}$.
    $F$ is a function on $I$ by
        \auto{free_module_regular_family_type,
            funs_is_function,funs}.
    For every $i\in I$ we have $F(i)$ is a left module over $R$ by
        \auto{free_module_regular_family_apply,
            left_regular_module_is_exact}.
    $g$ is a function on $I$ and for every $i\in I$ we have
        $g(i)$ is a linear map from $\LeftRegularModule{R}$ to $M$ over $R$ by
        \obv{free_module_coefficient_family_type}.
    Hence for every $i\in I$ we have $g(i)$ is a linear map from
        $F(i)$ to $M$ over $R$ by
        \auto{free_module_regular_family_apply}.
    $x\in\carrier[\DirectSumModule{R}{I}{F}]$ by
        \auto{free_module_constructor_eq}.
    $\apply{\directSumModuleMediator{R}{I}{F}{M}{g}}{x}
        =\finiteFold{\ModuleAdditiveGroup{M}}{
            \dependentSupport{I}{x}{z}}{
            \directSumGroupContribution{I}{A}{
                \ModuleAdditiveGroup{M}}{g}{x}}$ by
        \auto{direct_sum_module_mediator_apply}.
    $z=I\times\{\ringZero[R]\}$ by
        \obv{free_module_zero_section}.
    $x\in\funs{I}{\carrier[R]}$ by
        \auto{free_module_carrier,
            finitely_supported_functions_iff}.
    $\dependentSupport{I}{x}{I\times\{\ringZero[R]\}}
        =\supportOn{I}{x}{\ringZero[R]}$ by
        \auto{ring_zero_type,constant_family_support}.
    Follows by \auto{free_module_lift}.
\end{proof}

\begin{theorem}\label{free_module_lift_unique}
    Let $R$ be a ring.
    Let $M$ be a left module over $R$.
    Let $h\in\funs{I}{\carrier[M]}$.
    Suppose $k$ is a linear map from $\FreeModule{R}{I}$ to $M$ over $R$.
    Suppose for every $i\in I$ we have
        $k(\freeBasisVector{R}{I}{i})=h(i)$.
    Then $k=\freeModuleLift{R}{I}{M}{h}$.
\end{theorem}
\begin{proof}
    Let $F=I\times\{\LeftRegularModule{R}\}$ and
        $S=\FreeModule{R}{I}$.
    Let
        $g=\Apply{\graphOnDomainOp}{
            I,\fun{i:\Set}{
                \Apply{\graphOnOp}{
                    \carrier[R],\carrier[M],
                    \fun{a:\Set}{\scalarMul[M](a,h(i))}}}}$.
    $F$ is a function on $I$ by
        \auto{free_module_regular_family_type,
            funs_is_function,funs}.
    For every $i\in I$ we have $F(i)$ is a left module over $R$ by
        \auto{free_module_regular_family_apply,
            left_regular_module_is_exact}.
    $g$ is a function on $I$ and for every $i\in I$ we have
        $g(i)$ is a linear map from $\LeftRegularModule{R}$ to $M$ over $R$ by
        \obv{free_module_coefficient_family_type}.
    Hence for every $i\in I$ we have $g(i)$ is a linear map from
        $F(i)$ to $M$ over $R$ by
        \auto{free_module_regular_family_apply}.
    Show that for every $i\in I$ we have
        $k\circ\directSumModuleInclusion{R}{I}{F}{i}=g(i)$.
    \begin{subproof}
        Fix $i\in I$.
        Let $u=\directSumModuleInclusion{R}{I}{F}{i}$ and
            $L=\LeftRegularModule{R}$.
        $F(i)=L$ by \auto{free_module_regular_family_apply}.
        $u$ is a linear map from $L$ to $S$ over $R$ by
            \auto{free_module_constructor_eq,
                direct_sum_module_inclusion_linear}.
        Hence $u\in\funs{\carrier[L]}{\carrier[S]}$ by
            \obv{linear_map_type}.
        $k\in\funs{\carrier[S]}{\carrier[M]}$ by
            \auto{linear_map_type}.
        $k$ is composable with $u$ by \auto{funs_ran,funs}.
        Hence $k\circ u\in\funs{\carrier[L]}{\carrier[M]}$ by
            \obv{funs_circ}.
        $g(i)\in\funs{\carrier[L]}{\carrier[M]}$ by
            \auto{linear_map_type}.
        Show that for every $a\in\carrier[L]$ we have
            $\apply{k\circ u}{a}=\apply{g(i)}{a}$.
        \begin{subproof}
            Fix $a\in\carrier[L]$.
            Hence $a\in\carrier[R]$ by
                \auto{left_regular_module_carrier}.
            $\ringOne[R]\in\carrier[L]$ by
                \auto{left_regular_module_carrier,ring_one_type}.
            $\ringOne[R]\in\carrier[R]$ by \obv{ring_one_type}.
            $\scalarMul[L](a,\ringOne[R])
                =\ringMul[R](a,\ringOne[R])$ by
                \auto{left_regular_scalar_apply}.
            $\ringMul[R](a,\ringOne[R])=a$ by
                \obv{ring_mul_one_right}.
            Hence $\scalarMul[L](a,\ringOne[R])=a$.
            $u(a)=\scalarMul[S](a,u(\ringOne[R]))$ by
                \auto{linear_map_scalar}.
            $u(\ringOne[R])=\freeBasisVector{R}{I}{i}$ by
                \auto{free_basis_vector}.
            Hence $u(a)=\scalarMul[S](a,
                \freeBasisVector{R}{I}{i})$.
            $\freeBasisVector{R}{I}{i}\in\carrier[S]$ by
                \auto{free_basis_vector_type}.
            Hence
                $k(\scalarMul[S](a,\freeBasisVector{R}{I}{i}))
                    =\scalarMul[M](a,k(\freeBasisVector{R}{I}{i}))$ by
                \auto{linear_map_scalar}.
            $k(\freeBasisVector{R}{I}{i})=h(i)$.
            Hence $k(u(a))=\scalarMul[M](a,h(i))$.
            $k$ and $u$ are functions by
                \auto{funs_is_function}.
            $a\in\dom{u}$ by \auto{funs}.
            $\apply{k\circ u}{a}=k(u(a))$ by \obv{circ_apply}.
            $\apply{g(i)}{a}=\scalarMul[M](a,h(i))$ by
                \obv{free_module_coefficient_family_apply}.
        \end{subproof}
        Follows by \auto{typed_functions_equal_from_agreement}.
    \end{subproof}
    $k$ is a linear map from $\DirectSumModule{R}{I}{F}$ to $M$ over $R$ by
        \auto{free_module_constructor_eq}.
    Hence $k=\directSumModuleMediator{R}{I}{F}{M}{g}$ by
        \obv{direct_sum_module_mediator_unique}.
    Follows by \auto{free_module_lift}.
\end{proof}
